\chapter{Introducción}

\label{cap:introduccion}


\section{Antecedentes}

La medicina nuclear ha experimentado en las últimas décadas una transición hacia la medicina de precisión, impulsada en gran medida
por el auge de la \textbf{teranosis} (acrónimo de terapia y diagnóstico). Este enfoque se basa en el uso de pares de isótopos radiactivos
("matched pairs") que comparten propiedades químicas idénticas o muy similares, permitiendo utilizar uno para el diagnóstico por
imagen y otro para la terapia en la misma diana molecular.

Uno de los pares teranósticos más relevantes en la práctica clínica actual es el formado por el \textbf{Itrio-86} (\yochoseis) y el
\textbf{Itrio-90} (\ynuevecero). El \ynuevecero es un emisor beta negativo puro ($\beta^-$). Sin embargo, su emisión carece de
fotones gamma detectables, dificulta su detección mediante las técnicas de imagen convencionales, imposibilitando una dosimetría
precisa previa al tratamiento. Su isótopo análogo, el \yochoseis es un emisor de positrones ($\beta^+$) que permite visualizar la
biodistribución del fármaco mediante Tomografía por Emisión de Positrones (PET).

A pesar de su utilidad, la imagen cuantitativa con \yochoseis presenta desafíos físicos significativos. Se clasifica como un
\textbf{isótopo "sucio"}, ya que ya que su desintegración conlleva la emisión simultánea de una cascada de rayos gamma de alta
energía (\textit{prompt gammas}). La literatura científica (\cite{Lubberink2011,Conti2016}) señala que estos gammas generan
coincidencias triples y aumentan la tasa de eventos aleatorios, lo que los escáneres comerciales estándar malinterpretan,
degradando severamente el contraste y la exactitud cuantitativa (\cite{Buchholz2003}).

Para abordar esta problemática desde una perspectiva industrial y tecnológica, la simulación Monte Carlo es la herramienta
fundamental. En este trabajo se utiliza el framework \textbf{\penred} (\cite{Gimnez-Alventosa2021})
basado en PENELOPE, para modelar con alta fidelidad la geometría y la física de detección del escáner preclínico
\textbf{Bruker PET/CT Si78}. El uso de las herramientas de reconstrucción propietarias de la empresa (Bruker RecoServer)
permite validar los resultados en un entorno de aplicación real.

\section{Motivación}

La motivación principal de este trabajo surge de una necesidad clínica y tecnológica no resuelta: \textbf{la falta de precisión
cuantitativa} en las imágenes PET de \yochoseis impide aprovechar todo el potencial del par teranóstico \yochoseis/\ynuevecero.

